% Restored and adapted detailed content from the previous playground and verification draft.

% Source: paper/chapters/05-implementation/playground.tex

\section{The Demonstration Playground}\label{detail-sec:implementation-playground}

While the core contribution of this thesis is the compiler, a web-based \textit{Playground} was developed to demonstrate
the language's capabilities and provide an intuitive environment for composition.
The playground integrates the C++ compiler into a modern full-stack web application.

\subsection{System Architecture}\label{detail-subsec:system-architecture}

The playground follows a distributed architecture:

\noindent \textbf{Frontend (Next.js).} A React-based interface that provides the editor and visualization tools.

\noindent \textbf{Backend (Express).} A Node.js API that acts as a bridge.
It receives Motivo source code from the client, invokes the C++ \texttt{motivoc} binary via a sub-process, and streams
the resulting MIDI file back to the browser.

\noindent \textbf{Integration.} The system is containerized using Docker, ensuring that the C++ environment (including
Bison/Flex runtimes) is identical across development and production.

\subsection{IDE Integration and Diagnostics}\label{detail-subsec:ide-integration-and-diagnostics}

The playground features an IDE-like experience powered by the Monaco Editor.
A custom ``Monarch'' tokenizer was implemented for Motivo to provide syntax highlighting for musical keywords, pitch
literals, and control flow.

The most significant integration is the \textit{Diagnostics Feedback Loop}.
When a compilation fails, the Express server parses the structured error output from the compiler and returns it to the
client.
The React application then converts these diagnostics into Monaco ``Markers'', which underline the erroneous code and
display the compiler's semantic or syntax error messages in a tooltip.
This real-time feedback is enabled by the fast local compilation times discussed in the evaluation chapter.

\subsection{The Piano Roll and Synthesis}\label{detail-subsec:the-piano-roll-and-synthesis}

To visualize the composition, a custom \textit{Piano Roll} component was built using the HTML5 Canvas API\@.
Upon a successful compilation, the browser parses the returned MIDI bytes and renders them as colored rectangles on a
scrolling timeline.

The synthesis engine utilizes the \textit{Tone.js} library.
Because MIDI files contain only instructions and no sound, the playground loads a set of high-quality
\textit{Soundfonts} representing General MIDI instruments.
The playback logic schedules notes in the browser's audio context, allowing the user to hear their ``coded'' composition
instantly.
The synchronization between the scrolling piano roll and the audio playback is maintained through a high-precision
request-animation-frame loop, ensuring a professional and ``alive'' feel to the prototype.

% Source: paper/chapters/05-implementation/verification.tex

\section{Verification and Quality Assurance}\label{detail-sec:implementation-verification}

To ensure the correctness of the compilation pipeline and the validity of the generated MIDI artifacts, a comprehensive
testing suite was implemented using the \textit{Google Test} framework.
The testing strategy follows a tiered approach, covering individual units, module boundaries, and end-to-end
integration.

\subsection{Unit Testing}\label{detail-subsec:unit-testing}

Unit tests focus on the smallest verifiable components of the compiler.
This is particularly important for the musical domain logic, where small errors in calculation can lead to significant
musical dissonance.
Key test cases include:

\noindent \textbf{Pitch Math.} Verifying that pitch literals (e.g., \texttt{A4}) map correctly to MIDI note
numbers (e.g.,
69) and that octave offsets are handled according to the standard.

\noindent \textbf{Diagnostic Reporting.} Ensuring that the \texttt{DiagnosticsEngine} correctly captures and formats
errors across different stages of the pipeline.

\noindent \textbf{Lexer Accuracy.} Stress-testing the Flex regex patterns with edge cases, such as multiple consecutive
accidentals or varying whitespace.

\subsection{Integration and Regression Testing}\label{detail-subsec:integration-and-regression-testing}

Integration tests verify that the different stages of the compiler work together seamlessly.
These tests utilize a ``Golden File'' methodology:

\noindent A set of representative Motivo source files (located in \texttt{tests/data/}) is compiled.

\noindent The resulting IR is checked for temporal consistency (e.g., ensuring no notes overlap unless explicitly
requested).

\noindent The final MIDI binary is compared against a previously validated ``golden'' version to detect
unexpected changes
in behavior.

Regression tests were instrumental during the refactoring process (e.g., when decomposing the monolithic lowerer).
By maintaining a suite of complex compositions---such as the \texttt{fur\_elise.motivo} example---the compiler's
output could be
verified for byte-level consistency after every architectural change.

\subsection{Integration in the Development Lifecycle}\label{detail-subsec:integration-in-the-development-lifecycle}

The testing suite is integrated into the build system via CMake's \texttt{CTest} module.
This allows for one-command verification of the entire system (\texttt{make test}).
By enforcing a high level of test coverage, particularly in the lowering and backend stages, the implementation provides
high confidence that the generated compositions are mathematically and musically accurate.
This rigorous verification process is a cornerstone of the thesis, transforming a creative tool into a reliable software
system.
